% Julia 开发笔记
% license Xiao
% type Note

本文是关于 Julia 解释器的原理： 它使用了哪些技术， 可以使得它作为一门动态语言能达到编译语言的性能。

\begin{itemize}
\item \href{https://docs.julialang.org/en/v1/devdocs/init/}{Julia 开发者文档}
\item C 语言这样的叫做 \textbf{Ahead of Time（AOT）} 编译。
\item Julia 完全使用 \textbf{Just-In-Time（JIT）}。 并不是把部分代码使用该优化，而是全部。Julia 默认不直接理解并执行高级语言（一个解释器见\href{https://juliadebug.github.io/JuliaInterpreter.jl/stable/}{这里}）。 每个重载的函数第一次被调用的时候， julia 都会先给它生成专门的 llvm 机器码。 这个过程还有一个不那么正规的名字叫做 \textbf{Just Ahead-of-Time（JAOT）}。
\item \textbf{Multiple Dispatch （MD）}： 可以基于参数的类型生成不同的专门代码。
\item \textbf{Type Inference}： 自动推导变量类型。
\item 内建并行， 包括分布计算
\item 高效内存管理： 减少内存分配， 增加内存重复利用。
\item 基于 \textbf{LLVM}： Julia 编译器先把 Julia 语言变为 \textbf{Julia IR} （和 LLVM IR 相似，而且有若干层）， 再变为 \enref{LLVM IR}{llvmIR}, 最后交给 LLVM 进行优化。
\item Julia IR 后期使用 \textbf{Single Static Assignment (SSA)} 形式， 把 julia 代码的控制流表示为 \textbf{directed acyclic graph (DAG)} 的 CFG（Control Flow Graph）。 Julia IR 中的变量类型都是明确的。
\item \verb`LLVM.jl` 包可以把 julia 代码直接生成 LLVM IR， 或者把 Julia IR 转为 LLVM IR。
\item 一些预备知识： 编译原理， 编译器设计， LLVM， Julia 背后原理。
\item 关于\href{https://zhuanlan.zhihu.com/p/569704329}{多重分派}
\item \href{https://zhuanlan.zhihu.com/p/515863922}{Julia 的批评}
\end{itemize}

\begin{figure}[ht]
\centering
\includegraphics[width=14.2cm]{./figures/7293eeb0d6cd9c43.png}
\caption{Julia JIT 流程} \label{fig_julia0_1}
\end{figure}

\subsection{Julia 的 AST 格式}
\begin{itemize}
\item \verb`dump(表达式)` 函数可以显示语法树，如 \verb`expr = :(x + 2y)` 显示
\begin{lstlisting}[language=julia]
Expr
  head: Symbol call
  args: Array{Any}((3,))
    1: Symbol +
    2: Symbol x
    3: Expr
      head: Symbol call
      args: Array{Any}((3,))
        1: Symbol *
        2: Int64 2
        3: Symbol y
\end{lstlisting}
\item \verb`Expr` 类型是内建的，似乎无法查看源码。
\item \verb`dump(类型)` 可以显示该类型的定义。 \verb`dump(Expr)` 会打印 
\begin{lstlisting}[language=julia]
Expr <: Any
  head::Symbol
  args::Vector{Any}
\end{lstlisting}
\item 创建表达式如 \verb`ex = Expr(:call, :+, 1, 2)`
\item \verb`@which dump(:(x + 2y))` 可以看到该函数在哪里定义（\verb`C:\Users\addis\AppData\Local\Programs\Julia-1.11.7\share\julia\base\show.jl`）。
\item \verb`dump(arg; maxdepth=最大层数)` 可以指定最大层数。
\item 接下来调用 \verb`function dump(io::IOContext, @nospecialize(x), n::Int, indent)`
\item 但 parser 似乎不是用 Julia 写的，因为不如调试无法看到生成语法树的过程。
\item 插值表达式 \verb`$(...)` 可以出现在 \verb`quote ... end` 中，直接把一部分表达式求值并插入而不是把 \verb`...` 作为表达式插入。 例如 \verb`:(a + b)` 中 \verb`+, a, b` 都是 Symbol， 但 \verb`a = 1; ex = :($(a+1) + b)` 得到的就是 \verb`:(2 + b)`
\end{itemize}


\subsection{把 Julia 代码编译成 LLVM IR 代码}
首先安装 \verb`using Pkg; Pkg.add("LLVM")` （其实没必要）

在 REPL 中输入如下代码
\begin{lstlisting}[language=julia,caption=julia]
function add(x::Int, y::Int)::Int
    return x + y
end
lowered_ir = @code_lowered add(1, 2) # 生成 Lowered Julia IR
typed_ir = @code_typed add(1, 2) # 生成静态类型的 Julia IR
llvm_ir = @code_llvm add(1, 2) # 生成 LLVM IR
\end{lstlisting}

也可以把函数定义写到 jl 脚本中，然后在 REPL 中 \verb`include("脚本名.jl")`， 再 \verb`@code_xxx add(1, 2)`。 \verb`include()` 相当于把脚本代码直接插入到该位置。 如果找不到，就用 \verb`using InteractiveUtils; a = InteractiveUtils.@code_lowered add(1, 2);`

注意只会生成指定函数的 IR，不会递归生成它调用的其他函数的 IR。

\begin{lstlisting}[language=none,caption=julia ir]
julia> lowered_ir = @code_lowered add(1, 2)
CodeInfo(
1 ─ %1  = Main.Int
│   %2  = Main.:+
│   %3  = (%2)(x, y)   
│         @_4 = %3
│   %5  = @_4
│   %6  = %5 isa %1
└──       goto #3 if not %6
2 ─       goto #4
3 ─ %9  = @_4
│   %10 = Base.convert(%1, %9)
└──       @_4 = Core.typeassert(%10, %1)
4 ┄ %12 = @_4
└──       return %12
)
\end{lstlisting}

生成的 LLVM IR 代码如下
\begin{figure}[ht]
\centering
\includegraphics[width=14.25cm]{./figures/f34505db9ea28f8f.png}
\caption{LLVM IR} \label{fig_julia0_2}
\end{figure}

注意返回的 \verb`lowered_ir` 变量并不是一个字符串，而是结构化的信息。

\subsubsection{如何读懂 Julia IR}
如果 Julia IR 中的东西不懂，文档也不完善，那就直接执行试试即可， Julia IR 中的函数都是可以在 REPL 中执行的。 例如上面的 \verb`Main.Int` 执行出来就是 \verb`Int64::DataType`。 找许多 Julia 源码和 IR 对照下来基本就能明白。

\subsection{循环的 IR}
代码
\begin{lstlisting}[language=julia]
for n in 1:5
    c += n
end
\end{lstlisting}
对应的 IR 是
\begin{lstlisting}[language=none]
14 ┄ %77  = Main.:(:)                         ; for n in 1:5
                                              ;     % ==循环准备阶段==
│    %78  = (%77)(1, 5)                       ;     % 返回 ::UnitRange{Int64}
│           @_4 = Base.iterate(%78)           ;     % 初始化迭代器： (值, 循环指数)
                                              ;     % 类型 Tuple{Int64, Int64}
                                              ;     % ， 当前为 (1, 1)
│    %80  = @_4
│    %81  = %80 === nothing
│    %82  = Base.not_int(%81)                 ;     % 非运算
└───        goto #17 if not %82
                                              ;     % ==循环体==
15 ┄ %84  = @_4
│           n = Core.getfield(%84, 1)
│    %86  = Core.getfield(%84, 2)             ;     % 循环指数（对 UnitRange{Int64}
                                              ;     %    未必从 1 开始也未必步长为 1）
│    %87  = Main.:+                           ;     c += n
│    %88  = c
│    %89  = n
│           c = (%87)(%88, %89)
│           @_4 = Base.iterate(%78, %86)      ;     % 迭代（nothing 表示结束）
│    %92  = @_4
│    %93  = %92 === nothing
│    %94  = Base.not_int(%93)
└───        goto #17 if not %94
16 ─        goto #15
                                              ; end
\end{lstlisting}

代码
\begin{lstlisting}[language=julia]
function foo()
a = 0
for i in [1, "123", 1+2im]
    a += 1
    println(i, a);
end
end
\end{lstlisting}
对应的 IR 是
\begin{lstlisting}[language=none]
CodeInfo(
1 ─       a = 0
                                         ; % ==循环准备==
│   %2  = Main.:+                        ; % 构造 [1, "123", 1+2im]
│   %3  = Main.:*
│   %4  = Main.im
│   %5  = (%3)(2, %4)
│   %6  = (%2)(1, %5)
│   %7  = Base.vect(1, "123", %6)
│         @_2 = Base.iterate(%7)         ; % 迭代器
│   %9  = @_2
│   %10 = %9 === nothing
│   %11 = Base.not_int(%10)
└──       goto #4 if not %11
2 ┄ %13 = @_2                            ; % ==循环体==
│         i = Core.getfield(%13, 1)      ; % 循环变量
│   %15 = Core.getfield(%13, 2)          ; % 循环指数
│   %16 = Main.:+
│   %17 = a
│         a = (%16)(%17, 1)
│   %19 = i
│   %20 = a
│         Main.println(%19, %20)
│         @_2 = Base.iterate(%7, %15)
│   %23 = @_2
│   %24 = %23 === nothing
│   %25 = Base.not_int(%24)
└──       goto #4 if not %25
3 ─       goto #2
4 ┄       return nothing
)
\end{lstlisting}

\subsection{Core 模块}
\begin{itemize}
\item \verb`Core._apply_iterate(Base.iterate, 函数, 容器1, 容器2, ...)` 会依次把每个 \verb`容器` 的每个元素作为入参去调用 \verb`函数`。 其中 \verb`Base.iterate` 是一个函数，用途未知。
\end{itemize}

\subsection{Base 模块}
\begin{itemize}
\item \verb`Base.not_int(变量)` 逻辑非运算（似乎非逻辑类型就按 bit 取非）
\item \verb`Base.Pairs(("a","b"), (1,2))` 返回的类型是 \verb`Base.Pairs{Int64, String, Tuple{Int64, Int64}, Tuple{String, String}}`，但类似于 \verb`Dict([1=>"a", 2=>"b"])`。 可以使用 \verb`var[key]` 获取对应的值。
\item \verb`Base.iterate(容器)` 可以得到 \verb`容器` 的初始迭代器。 \verb`Base.iterate(容器, 迭代器)` 可以得到下一个元素的迭代器。
\end{itemize}

\subsection{Main 模块}

\begin{itemize}
\item 基本所有用户能调用的常用内建函数和算符都是 \verb`Main` 模块的， 例如 \verb`Main.println`， \verb`Main.error`， \verb`Main.sqrt` 等。
\item 算符表示为函数 \verb`Main.:算符符号`， 如大于号 \verb`Main.:>`， 减号 \verb`Main.:-`， 不等于 \verb`Main.:!=`。 \verb`1:5` 中的冒号表示为 \verb`Main.:(:)`
\end{itemize}

\subsection{code_lowered 源码分析}
\begin{itemize}
\item \verb`@code_lowered` 会调用 \verb`C:\Users\addis\AppData\Local\Programs\Julia-1.11.6\share\julia\base\reflection.jl` 中的 \verb`function code_lowered(@nospecialize(f), @nospecialize(t=Tuple); generated::Bool=true, debuginfo::Symbol=:default)`
\item \verb`reflection.jl` 的 \verb`function method_instances(@nospecialize(f), @nospecialize(t), world::UInt)` 会返回 \verb`Core.MethodInstance` 数组，每个元素是 \verb`Core.Compiler.specialize_method(match)` 返回的。 其 \verb`source` 字段应该就是压缩后的 IR
\item \verb`reflection.jl` 的 \verb`uncompressed_ir(m::Method)` 最后会调用 \verb`ccall(:jl_uncompress_ir, Any, (Any, Ptr{Cvoid}, Any), m, C_NULL, s)::CodeInfo`，也就是调用 c 语言函数。
\item 所以看来 code_lowered 也是 C/C++ 实现的
\item \verb`@code_lowered` 返回一个 \verb`Core.CodeInfo` 类型的结构体，最后经过一次 \verb`remove_linenums!(code)`，打印出来就是我们看到的文本格式 Julia IR，其 \verb`code` 字段（\verb`Vector{Any}`）就是内部的数据结构。
\item \verb`%数字` 并不是连续的，而是行号也就是 \verb`code[数字]`
\item 每个元素可能是 \verb`Expr`, \verb`GlobalRef`（例如 \verb`Main.:+`）, \verb`Core.SlotNumber`（变量名赋值给临时变量）, \verb`Core.NewVarNode`, \verb`Core.GotoIfNot`（goto 语句）
\item 变量名在 \verb`code` 中用符号 \verb`:_数字` 来表示，具体的变量名对照表在 \verb`slotnames` 字段中 \verb`slotnames[数字]`。
\item \verb`fieldnames(GlobalRef)` 为 \verb`(:mod, :name, :binding)`， 是为了引用某个模块中的某个全局对象。 \verb`mod::Module` 字段是模块名，\verb`name::Symbol`，是被引用的符号名， \verb`binding::Core.Binding` 未知。
\end{itemize}

\subsection{code_typed} 源码分析
\begin{itemize}
\item 返回的第一个出参是 \verb`Core.CodeInfo` 类型的
\item 之前 untyped 没有常量折叠，typed 有，以及死分支剪切
\item \verb`code` 字段是 \verb`Vector{Any}`， 元素类型有 \verb`Expr`, \verb`Nothing`, \verb`Core.ReturnNode`（表示 return）。其中 \verb`Expr` 的 \verb`head` 一般是 \verb`:invoke`， \verb`args` 分别是被调函数的 \verb`MethodInstance`， 和被调函数的 \verb`GlobalRef`（例如 \verb`Main.print`），函数入参1，入参2，……
\item \verb`Core.ReturnNode` 只有一个 \verb`val` 字段，就是返回值。
\item phi 节点表示为 \verb`Core.PhiNode`
\item 还有一个奇怪的 pi 节点 \verb`Core.PiNode`，有两个字段 \verb`val, typ`
\end{itemize}
